% Front matter: abstract, keywords, RQ, hypotheses, at-a-glance. Written LAST per plan.

\begin{abstract}
Retailers live or die by two questions: how much will customers want tomorrow,
and how much stock should sit on the shelf waiting for them? Most forecasting
research answers only the first question. This study answers both. We compare
twelve forecasting approaches --- from the humble moving average to a global
long short-term memory (LSTM) neural network --- on two retail demand
environments: 500 intermittent-demand series from the Walmart M5 competition
data and all 500 series of a dense store-item demand dataset. Every model faces
identical conditions: the same calendar window, the same 28-day horizon over
eight rolling origins, the same error metrics, and --- crucially --- the same
downstream inventory policy, so that forecasts are finally judged by the costs
they create, not just the errors they contain.

The headline result is a tale of two datasets. On sparse M5 demand, the LSTM is
the best forecaster (mean absolute scaled error, MASE, of 1.316) and also the
cheapest inventory driver (cost 152.83), holding first place across 25 of 27
tested inventory policies: a robust conclusion. On dense store demand, the LSTM
is again the best forecaster (MASE 0.978) but only the fourth-cheapest inventory
driver (cost 2,247.45), beaten by a simple moving average (2,084.49); leadership
there changes hands with the policy (moving average 9, SES 7, naive 6, DES 5 of
27 policies): a fragile, policy-sensitive conclusion. In plain terms, forecast
accuracy does not automatically buy better shelf decisions --- sometimes the
smartest predictor is not the cheapest stocker. Ninety-one paired statistical
comparisons (Holm-corrected) separate real differences from noise, and a
27-policy sensitivity grid separates conclusions an executive can bank on from
ones that flip when costs or lead times move. All data, code, and validation
gates are frozen and traceable; the large-language-model phase of this research
program is designed but deliberately unexecuted and appears here only as future
work.
\end{abstract}

\noindent\textbf{Keywords:} retail demand forecasting; inventory optimization;
intermittent demand; M5 competition; LSTM; Croston method; forecast evaluation;
MASE; order-up-to policy; sensitivity analysis.

\subsection*{Research question}
How does the effectiveness of AI-based inventory optimization change as we move
from traditional forecasting models toward neural-network approaches, when
effectiveness is measured at the level of inventory outcomes rather than
forecast errors alone?

\subsection*{Hypotheses (neutral framing)}
\begin{description}[leftmargin=!,labelwidth=2.2cm]
\item[H1.] Forecast accuracy differs significantly across model generations on
the same series and horizons.
\item[H2.] Forecast-accuracy differences translate into measurable
inventory-outcome differences under a common policy.
\item[H3.] Increasing model sophistication yields monotonically improving
inventory outcomes. (The claim under test --- rejection is a publishable
finding, not a failure.)
\item[H4.] Observed patterns are consistent across the two retail environments.
\end{description}

\subsection*{At a glance}
\begin{center}
\begin{tabular}{p{3.2cm}p{5.4cm}p{5.4cm}}
\toprule
 & \textbf{M5 (sparse, intermittent)} & \textbf{Store (dense, smooth)} \\
\midrule
Best forecaster & LSTM (MASE 1.316) & LSTM (MASE 0.978) \\
Best inventory driver & LSTM (cost 152.83) & Moving average (cost 2,084.49) \\
Verdict & \textbf{Robust:} LSTM first in 25/27 policies & \textbf{Fragile:} lead shared 9 / 7 / 6 / 5 \\
Lesson & Best forecast $\rightarrow$ best decision & Best forecast $\not\rightarrow$ best decision \\
\bottomrule
\end{tabular}
\end{center}
\begin{center}
\emph{Forecast accuracy does not necessarily imply inventory decision superiority.}
\end{center}
